\section{Conclusions and Future Work}
This paper has presented a lightweight TypeScript-based 3D engine for browser-delivered interactive educational content, together with three targeted optimization strategies addressing distinct dimensions of runtime performance. The engine wraps a low-level graphics library behind a high-level component model inspired by mature commercial engines, enabling course content teams to build interactive 3D learning environments using familiar software-engineering patterns while never directly encountering the underlying graphics API.

The empirical evaluation demonstrates that the proposed engine sustains interactive frame rates exceeding 100 FPS on a mid-range gaming laptop and a stable 60 FPS on a desktop workstation. The full optimization stack produces a 60.1\% reduction in maximum frame time and a 97.5\% reduction in frame-time variability on the laptop tier, while imposing no measurable cost on the desktop tier. The optimizations therefore correctly target the lower-capability portion of the hardware spectrum that institutional deployments cannot avoid. These results indicate that practical, targeted improvements at the engine architecture and asset pipeline levels can substantially improve the perceived responsiveness of browser-based educational 3D applications, without requiring the heavyweight machinery of a full native engine runtime exported to WebAssembly.

Future work includes extending the system to a broader range of devices, including integrated graphics and mobile platforms, as well as introducing a \texttt{WebGPU} backend to further improve rendering efficiency. Additional planned extensions include geometry batching, level-of-detail (LOD) systems, and a graphical scenario authoring tool designed for educators without programming expertise.